\documentclass[11pt,a4paper]{article}
\usepackage[utf8]{inputenc}
\usepackage[T1]{fontenc}
\usepackage[french]{babel}
\usepackage{geometry}
\geometry{margin=2.5cm}
\usepackage{setspace}
\onehalfspacing
\usepackage{amsmath, amssymb}
\usepackage{graphicx}
\usepackage{float}
\usepackage{caption}
\usepackage[hidelinks]{hyperref}
\usepackage{microtype}
\usepackage{textcomp}

\newcommand{\LABeCOdeux}{LABeCO\textsubscript{2}}

% -------------------------------------------------
% Début du document
% -------------------------------------------------
\begin{document}

\section{Mesure et standardisation de l’empreinte carbone des pratiques de recherche académique}

\paragraph{Positionnement général}

Ce travail vise à aborder la question de l’empreinte carbone de la recherche non pas comme un indicateur seulement informatif ou bien punitif,
mais comme un \textbf{problème scientifique de mesure, de définition et de standardisation}.
Dans la logique de mon projet de recherche, il vise à interroger les pratiques expérimentales des chercheurs\textperiodcentered euses, avec leurs contraintes matérielles, et la manière dont ces pratiques peuvent être quantifiées et
comparées de façon rigoureuse sur la base de leur coût environnemental, permettant ainsi un regard critique sur nos comportements et de nos protocoles de recherche.

L’enjeu central n’est pas, là encore, la performance du calcul en lui-même, mais la \textbf{pertinence de ce que l'on mesure}, la robustesse du cadre méthodologique, et l'adaptation des outils développés
aux usages quotidiens des chercheurs\textperiodcentered euses.

\paragraph{LABeCO\textsubscript{2} : un logiciel comme instrument scientifique}

Ce travail s’appuie sur le développement de mon projet \textbf{LABeCO\textsubscript{2}}, un logiciel dédié à la quantification de l’empreinte carbone des activités de recherche individuelle (ou d'équipe),
dont la structure générale est déjà établie (et présentée dans mon expérience de recherche).
\textit{LABeCO\textsubscript{2}} n’est pas conçu comme un outil à l'échelle institutionnelle, mais comme un \textbf{instrument de recherche} individuel,
au même titre qu’un dispositif expérimental ou un outil d’analyse de données.

Un élément essentiel dans cette partie du projet est l'établissement d'un processus de réflexion sur le \textbf{format de l’outil}.
Contrairement à de nombreuses approches existantes, LABeCO\textsubscript{2} est pensé comme une application pure, exécutable en local, et non comme une application web.
Ce choix est motivé par les usages réels des chercheurs\textperiodcentered euses :
le calcul de l’empreinte carbone doit pouvoir s’intégrer de manière fluide aux pratiques quotidiennes (préparation d’expériences, rédaction, gestion de projets),
sans dépendance permanente à une connexion externe ni à des interfaces lourdes.

Le développement du logiciel constitue ainsi un objet scientifique à part entière :
\begin{itemize}
    \item conception d’une architecture adaptée à des usages récurrents et évolutifs,
    \item structuration des données et des résultats émis pour permettre des comparaisons entre protocoles,
    \item calcul des incertitudes et explicitation des hypothèses de calcul,
    \item possibilité d’évolution et d’appropriation par d’autres équipes.
\end{itemize}

\paragraph{Standardisation et définition du cadre de la recherche}

Au-delà de l’outil, un enjeu scientifique majeur de cet axe concerne la \textbf{définition du périmètre} de calcul de l’empreinte carbone d'une activité de recherche.
Contrairement à des approches globales ou institutionnelles, ce projet se concentre volontairement à l’échelle individuelle ou d'un projet expérimental spécifique.

Une question centrale est alors : \textbf{qu’est-ce qui relève du cadre de la recherche que l’on souhaite mesurer, et qu’est-ce qui n’en relève pas ?}
Cette réflexion implique de poser des frontières claires entre :
\begin{itemize}
    \item les activités directement liées à la production scientifique (instrumentation, consommables, campagnes de terrain, calculs, déplacements professionnels),
    \item les activités périphériques ou personnelles,
    \item les infrastructures mutualisées dont l’impact doit être traité de manière spécifique.
\end{itemize}

Le cadre n’est pas évident à mettre en place et se pose comme un sujet de recherche à part entière. Elle engage des choix méthodologiques,
des hypothèses explicites et une réflexion critique sur ce que l’on cherche réellement mesurer lorsque l’on parle d’empreinte carbone de la recherche.

Cette réflexion a d’ores et déjà été engagée au cours de présentations informelles et d’échanges avec de nombreux chercheurs\textperiodcentered euses issus de disciplines variées.
Ces discussions ont mis en évidence un fort intérêt pour la démarche, ainsi qu’un besoin partagé de disposer d’outils et de cadres clairs pour quantifier l'empreinte carbone à une échelle plus individuelle, et pas seulement institutionnelle.

Ces premiers retours me rendent confiant quant à la pertinence et à la faisabilité de ce volet, et m’ont permis d’identifier un premier noyau de chercheurs\textperiodcentered euses intéressé\textperiodcentered es à la formation d'un collectif pour
structurer le projet. Je souhaite également organiser rapidement un séminaire de réflexion.
Son objectif sera de structurer collectivement une première définition du périmètre de la recherche, de discuter des hypothèses de travail,
et d’élaborer progressivement un cadre partagé des données et des besoins.

\paragraph{Une démarche collaborative dans sa construction}

Dans la continuité de cette approche, ce cadre scientifique est conçu dès l’origine comme une \textbf{démarche collaborative}.
La standardisation du calcul de l’empreinte carbone de la recherche ne peut en effet pas reposer sur une définition figée et imposée,
mais nécessite une construction collective et évolutive, impliquant les chercheurs\textperiodcentered es eux-mêmes, selon leurs pratiques réelles.

Cette dimension collaborative intervient à deux niveaux complémentaires.
D’une part, dans la \textbf{définition du domaine de validité du calcul} : identifier ce qui relève effectivement du cadre de la recherche que l'on souhaite mesurer,
clarifier les frontières entre activités scientifiques, infrastructures mutualisées et autres activités non liées directement.
Ces choix méthodologiques conditionnent directement la pertinence des bilans carbone produits et constituent ainsi en eux-mêmes un objet de débat scientifique.

D’autre part, cette démarche collaborative est indispensable pour la \textbf{construction et l’enrichissement des bases de données} nécessaires au calcul.
Un seul acteur ne peut pas prendre en charge l'ensemble du calcul du coût environnemental de tous les objets de recherche.
L’amélioration de la précision des facteurs d’émission, en particulier pour les consommables, réactifs et dispositifs expérimentaux spécifiques,
ne peut se faire qu’à travers une contribution d'une grande part de la communauté scientifique.
Dans cette perspective, en marge de \textit{LABeCO\textsubscript{2}}, un second outil évolutif open-source sera ouvert : une base de données compléte des empreinte carbonnes des objets de recherches (consommables, informatiques, ...)
calculé à partir des nouveaux standards définis collectivement. 

Ce travail collaboratif vise ainsi à faire émerger, progressivement, une \textbf{pratique normalisée de la mesure du coût carbone des activités de recherche}, basée sur une réflexion collective,
qui s'ancre dans les préoccupations modernes sur la sobriété.

\paragraph{Vers une méthodologie reproductible et comparable}

L’objectif de ce travail est de proposer une méthodologie claire, documentée et reproductible, permettant :
\begin{itemize}
    \item de comparer différentes stratégies expérimentales produisant une information scientifique équivalente,
    \item d’identifier les postes dominants d’impact carbone dans une activité de recherche,
    \item de rendre explicites les compromis entre frugalité instrumentale, robustesse et contraintes expérimentales.
\end{itemize} 

Cette standardisation ne vise pas à produire une valeur définitive, mais à établir un cadre commun de calcul, rendant les résultats comparables entre projets.
Cela permet l'ouverture à un regard critique et objectivé sur nos pratiques de recherche, permettant d'identifier d'éventuelles marges d'adaptation afin de réduire l'impact environnemental global de la recherche
sans en compromettre la performance (fig. \ref{fig:projet_phase5_Empreinte_carbone}). 

\paragraph{Portée scientifique}

En intégrant la mesure de l’empreinte carbone comme un paramètre explicite des pratiques expérimentales, ce cadre de recherche prolonge naturellement les objectifs du projet :
il s’agit de replacer les contraintes matérielles, énergétiques et instrumentales au coeur de la réflexion scientifique.

J’intégrerai, au sein de chacun de mes projets, le calcul de l’empreinte carbone et la définition explicite du périmètre de la recherche associée.
Ces premiers cadres auront vocation à servir de base à la discussion, à la remise en question et à l’évolution de ce périmètre.
Ils me permettront également d’identifier rapidement les principaux leviers sur lesquels agir afin de limiter, lorsque cela est possible, l’impact environnemental de mes activités de recherche.

\begin{figure}[H]
    \centering
    \includegraphics[width=0.9\textwidth]{/Users/souchaud/Documents/Travail/Transition_ecologique/LABeCO2/Formalisme/images/empreinte carbone.pdf}
    \captionsetup{font=small}
    \caption{
        Schéma conceptuel du volet transversal de mesure de l’empreinte carbone des pratiques de recherche.
        Le schéma illustre l’articulation entre les pratiques expérimentales de recherche, l’outil LABeCO\textsubscript{2} en tant qu’instrument de mesure, et la construction d’indicateurs d’empreinte carbone reproductibles et comparables.
        La démarche repose sur une construction collaborative impliquant les chercheur \textperiodcentered euses à chaque étape, depuis la définition du périmètre de la recherche jusqu’à l’enrichissement des bases de données et l’identification de leviers de réduction
        de l'emprunte carbone des pratiques de la recheche.   
            }
    \label{fig:projet_phase5_Empreinte_carbone}
\end{figure}


\end{document}